\begin{resumo}[Abstract]
\begin{otherlanguage*}{english}

This study investigates the fiscal sustainability of Brazilian state governments from 2002 to 2024, focusing on the response of the primary balance to public debt and on the role of debt refinancing contracts established under Federal Law No.~9,496/1997. The research adapts the empirical strategy of Abubakar, McCausland and Theodossiou (2025) to the Brazilian subnational context, using fiscal indicators as ratios to Current Net Revenue, given its institutional relevance in the Fiscal Responsibility Law and in state debt contracts. The first model estimates a fiscal reaction function in which the primary balance-to-revenue ratio responds to the lagged net consolidated debt-to-revenue ratio and to its interaction with the contractual revenue commitment ceiling. The model is estimated by two-stage least squares with state and year fixed effects, treating lagged debt and its interaction with the contractual ceiling as potentially endogenous variables. The second model examines the dynamics of state debt, considering debt persistence, the primary balance, economic growth and additional dynamic panel robustness specifications. The estimates aim to assess whether there is evidence consistent with Bohn's fiscal reaction approach and whether this reaction varies according to the permissiveness of contractual ceilings. The contribution of the study lies in shifting the focus from the mere existence of fiscal rules to their institutional design, exploring the heterogeneity of refinancing contracts as a relevant element for subnational fiscal sustainability.

\vspace{\onelineskip}

\noindent
\textbf{Keywords}: fiscal sustainability; state debt; fiscal rules; Law No.~9,496/1997; fiscal reaction function; Current Net Revenue.

\end{otherlanguage*}
\end{resumo}